\documentclass[14pt,a4paper]{extreport}
% preamble.tex — спільна преамбула для всіх лабораторних робіт. Підключається одразу після \documentclass у кожному lab*/main.tex.
%
% ЗБІРКА: LuaLaTeX (не pdfLaTeX). Причина: пакет listings під pdfLaTeX не читає UTF-8 і зупиняє збірку на кириличних коментарях у лістингах програмного коду. LuaLaTeX обробляє UTF-8 нативно.
%
% ШРИФТ: Liberation Serif — метрично сумісний з Times New Roman і містить повний набір кириличних гліфів (TeX Gyre Termes кирилиці не має, через що український текст не відображався взагалі).

\usepackage{fontspec}
\setmainfont{Liberation Serif}
\setmonofont{DejaVu Sans Mono}[Scale=0.82]

\usepackage{polyglossia}
\setmainlanguage{ukrainian}
\setotherlanguage{english}

\usepackage{microtype}

% Поля, інтервал, відступи й заголовки нижче взяті з шаблону НУЛП, за яким написана бакалаврська робота (../bachelor-thesis/thesis/main.tex). Той самий формат приймає В.М. Шиманський і для магістерських, тож усі звіти тримаємо в ньому, а не в тому, що переказували на консультації.
\usepackage[left=25mm,right=15mm,top=15mm,bottom=23mm,footskip=1.5mm]{geometry}
\usepackage{setspace}
% Саме \setstretch{1.5}, а не \onehalfspacing: останній дає близько 1.25.
\setstretch{1.5}

\usepackage{graphicx}
\usepackage{float}
\usepackage{booktabs}
\usepackage{array}
\usepackage{tabularx}
\usepackage{longtable}
\usepackage{enumitem}
% labelindent вирівнює маркер списку з абзацним відступом — так у шаблоні НУЛП.
\setlist{nosep}
\setlist[itemize]{label=\textbullet,leftmargin=*,labelindent=1.27cm}
\setlist[enumerate]{leftmargin=*,labelindent=1.27cm}
\usepackage{indentfirst}
\usepackage{amsmath}

\usepackage{hyperref}
\hypersetup{colorlinks=true,linkcolor=black,urlcolor=blue,citecolor=black}

% xurl дозволяє \url розривати адресу в будь-якому місці, а не лише після скісної риски. Без нього довгий сегмент шляху (напр. «scopus---podrobnaya- instruktsiya-dlya-uchenogo») виходить за межі сторінки, бо дефіс для hyperref не є точкою розриву. Вантажиться після hyperref — перевизначає його \UrlBreaks.
\usepackage{xurl}

% ---------------------------- Список використаної літератури ----------------------------
% biblatex + biber зі стилем IEEE: посилання «[1]» у тексті, записи — за порядком першої згадки. Обрано IEEE, а не gost-numeric (ДСТУ 8302:2015), — так вирішив автор звітів після порівняння обох варіантів.
% Кожен звіт із джерелами сам підключає свій references.bib через \addbibresource і друкує список \printbibliography[heading=none] під власним \section. Звіт без джерел від завантаженого biblatex нічого не втрачає.
% Оточення списку — стандартне для стилю: мітка від лівого поля, текст опису вирівняний після неї. Не підганяємо його під абзацний відступ нумерованих переліків тексту (labelindent у \setlist вище) — список літератури є окремим розділом зі своїм стандартним оформленням.
\usepackage[backend=biber,style=ieee]{biblatex}
% biblatex за замовчуванням додає розтяжний проміжок після «/ . : ? =» в адресах, і з вимкненими переносами рядок вирівнюється розрідженням самої адреси («https : / / www . scopus»). Рвати адресу й так дозволяє xurl, тож проміжок лише шкодить.
\setlength{\biburlbigskip}{0mu}
% Без додаткового інтервалу між описами: розділ має йти суцільним текстом, як решта звіту. Так само в шаблоні бакалаврської.
\setlength{\bibitemsep}{0pt}
\setlength{\bibparsep}{0pt}

% ---------------------------- Лістинги програмного коду ---------------------------- Використано minted (через Pygments), а не listings: listings не вміє коректно обробляти багатобайтові символи й «висмикує» кириличні коментарі з коду, зліплюючи їх в один рядок. Потребує прапорця -shell-escape.
\usepackage{xcolor}
\usepackage[newfloat]{minted}
% labelsep=period дає «Рис. 2.1. Назва» замість «Рис. 2: Назва» — формат НУЛП.
\usepackage[labelsep=period]{caption}

% baselinestretch=1 обов'язковий: без нього лістинг успадковує півторашку основного тексту, і довгий фрагмент коду перестає вміщатися у сторінку. Усередині [H]-флоата це не помилка збірки, а нескінченний цикл порожніх сторінок — lab1 так намолотив 161 тисячу сторінок.
\setminted{
  baselinestretch=1,
  fontsize=\footnotesize,
  linenos=true,
  frame=single,
  framesep=4pt,
  breaklines=true,
  breakanywhere=true,
  tabsize=4,
  numbersep=7pt,
  bgcolor=gray!8,
}
\setminted[python]{python3=true}

\SetupFloatingEnvironment{listing}{name=Лістинг,listname=Перелік лістингів}

% ---------------------------- Нумерація та оформлення заголовків ----------------------------
\usepackage{titlesec}

\setcounter{secnumdepth}{3}
\renewcommand{\thesection}{\arabic{section}}
\renewcommand{\thesubsection}{\arabic{section}.\arabic{subsection}}
\renewcommand{\thesubsubsection}{\arabic{section}.\arabic{subsection}.\arabic{subsubsection}}

% Розділ по центру, 16 пт, великими літерами; слово «РОЗДІЛ» опущене, бо звіт не дисертація. Номер лишається: \numberwithin прив'язує до нього нумерацію рисунків («Рис. 2.1»), і без нумерованих розділів вона б зламалася. \MakeUppercase четвертим аргументом отримує текст заголовка, тому в джерелі заголовки пишуться звичайним регістром.
\titleformat{\section}[block]
  {\centering\bfseries\fontsize{16}{19.2}\selectfont}
  {\thesection.}{0.3em}{\MakeUppercase}
% Підрозділи — зліва з абзацного відступу, звичайним регістром.
\titleformat{\subsection}[block]{\bfseries}{\thesubsection.}{1em}{}
\titleformat{\subsubsection}[block]{\bfseries\itshape}{\thesubsubsection.}{1em}{}

% Незіркові \titlespacing — обов'язкові. Без них titlesec застосовує зірковий варіант, який прибирає абзацний відступ у першого абзацу після заголовка й перекриває indentfirst. Значення — типові для report.
\titlespacing{\section}{0pt}{20pt}{6pt}
\titlespacing{\subsection}{1.27cm}{6pt}{0pt}
\titlespacing{\subsubsection}{1.27cm}{6pt}{0pt}

\setlength{\parindent}{1.27cm}
\setlength{\parskip}{0pt}
\renewcommand{\arraystretch}{1.3}

% Переноси слів вимкнені — так у шаблоні НУЛП. Натомість абзац розганяється міжслівними пробілами, тому tolerance і emergencystretch підняті: без них рядки виходили б за праве поле.
\hyphenpenalty=10000
\exhyphenpenalty=10000
\tolerance=9999
\setlength{\emergencystretch}{3em}

% ---------------------------- Підписи до рисунків і таблиць ----------------------------
\captionsetup[figure]{name=Рис.,justification=centering,font={bf,it}}
\captionsetup[table]{name=Таблиця,justification=raggedleft,singlelinecheck=false,labelfont=bf,textfont=normalfont}

% Рисунки, таблиці й формули нумеруються в межах розділу: «Рис. 2.1», а не «Рис. 5». Так у шаблоні НУЛП, і так само виглядали приклади на консультації.
\numberwithin{equation}{section}
\numberwithin{figure}{section}
\numberwithin{table}{section}

% ---------------------------- Колонтитули ----------------------------
\usepackage{fancyhdr}
\pagestyle{fancy}
\fancyhf{}
\fancyfoot[R]{\thepage}
\renewcommand{\headrulewidth}{0pt}
\fancypagestyle{plain}{%
  \fancyhf{}
  \fancyfoot[R]{\thepage}
  \renewcommand{\headrulewidth}{0pt}
}

% titlepage.tex — титульний аркуш, спільний для всіх дисциплін.
% Потребує визначених у subject.tex команд:
%   \SubjectName, \LecturerName, \StudentName, \StudentGroup, \AcademicYear
%
% Використання у тексті: \MakeLabTitle{1}{Тема роботи}
%
% Рід роботи. Дисципліни з практичними роботами перевизначають його після
% % preamble.tex — спільна преамбула для всіх лабораторних робіт. Підключається одразу після \documentclass у кожному lab*/main.tex.
%
% ЗБІРКА: LuaLaTeX (не pdfLaTeX). Причина: пакет listings під pdfLaTeX не читає UTF-8 і зупиняє збірку на кириличних коментарях у лістингах програмного коду. LuaLaTeX обробляє UTF-8 нативно.
%
% ШРИФТ: Liberation Serif — метрично сумісний з Times New Roman і містить повний набір кириличних гліфів (TeX Gyre Termes кирилиці не має, через що український текст не відображався взагалі).

\usepackage{fontspec}
\setmainfont{Liberation Serif}
\setmonofont{DejaVu Sans Mono}[Scale=0.82]

\usepackage{polyglossia}
\setmainlanguage{ukrainian}
\setotherlanguage{english}

\usepackage{microtype}

% Поля, інтервал, відступи й заголовки нижче взяті з шаблону НУЛП, за яким написана бакалаврська робота (../bachelor-thesis/thesis/main.tex). Той самий формат приймає В.М. Шиманський і для магістерських, тож усі звіти тримаємо в ньому, а не в тому, що переказували на консультації.
\usepackage[left=25mm,right=15mm,top=15mm,bottom=23mm,footskip=1.5mm]{geometry}
\usepackage{setspace}
% Саме \setstretch{1.5}, а не \onehalfspacing: останній дає близько 1.25.
\setstretch{1.5}

\usepackage{graphicx}
\usepackage{float}
\usepackage{booktabs}
\usepackage{array}
\usepackage{tabularx}
\usepackage{longtable}
\usepackage{enumitem}
% labelindent вирівнює маркер списку з абзацним відступом — так у шаблоні НУЛП.
\setlist{nosep}
\setlist[itemize]{label=\textbullet,leftmargin=*,labelindent=1.27cm}
\setlist[enumerate]{leftmargin=*,labelindent=1.27cm}
\usepackage{indentfirst}
\usepackage{amsmath}

\usepackage{hyperref}
\hypersetup{colorlinks=true,linkcolor=black,urlcolor=blue,citecolor=black}

% xurl дозволяє \url розривати адресу в будь-якому місці, а не лише після скісної риски. Без нього довгий сегмент шляху (напр. «scopus---podrobnaya- instruktsiya-dlya-uchenogo») виходить за межі сторінки, бо дефіс для hyperref не є точкою розриву. Вантажиться після hyperref — перевизначає його \UrlBreaks.
\usepackage{xurl}

% ---------------------------- Список використаної літератури ----------------------------
% biblatex + biber зі стилем IEEE: посилання «[1]» у тексті, записи — за порядком першої згадки. Обрано IEEE, а не gost-numeric (ДСТУ 8302:2015), — так вирішив автор звітів після порівняння обох варіантів.
% Кожен звіт із джерелами сам підключає свій references.bib через \addbibresource і друкує список \printbibliography[heading=none] під власним \section. Звіт без джерел від завантаженого biblatex нічого не втрачає.
% Оточення списку — стандартне для стилю: мітка від лівого поля, текст опису вирівняний після неї. Не підганяємо його під абзацний відступ нумерованих переліків тексту (labelindent у \setlist вище) — список літератури є окремим розділом зі своїм стандартним оформленням.
\usepackage[backend=biber,style=ieee]{biblatex}
% biblatex за замовчуванням додає розтяжний проміжок після «/ . : ? =» в адресах, і з вимкненими переносами рядок вирівнюється розрідженням самої адреси («https : / / www . scopus»). Рвати адресу й так дозволяє xurl, тож проміжок лише шкодить.
\setlength{\biburlbigskip}{0mu}
% Без додаткового інтервалу між описами: розділ має йти суцільним текстом, як решта звіту. Так само в шаблоні бакалаврської.
\setlength{\bibitemsep}{0pt}
\setlength{\bibparsep}{0pt}

% ---------------------------- Лістинги програмного коду ---------------------------- Використано minted (через Pygments), а не listings: listings не вміє коректно обробляти багатобайтові символи й «висмикує» кириличні коментарі з коду, зліплюючи їх в один рядок. Потребує прапорця -shell-escape.
\usepackage{xcolor}
\usepackage[newfloat]{minted}
% labelsep=period дає «Рис. 2.1. Назва» замість «Рис. 2: Назва» — формат НУЛП.
\usepackage[labelsep=period]{caption}

% baselinestretch=1 обов'язковий: без нього лістинг успадковує півторашку основного тексту, і довгий фрагмент коду перестає вміщатися у сторінку. Усередині [H]-флоата це не помилка збірки, а нескінченний цикл порожніх сторінок — lab1 так намолотив 161 тисячу сторінок.
\setminted{
  baselinestretch=1,
  fontsize=\footnotesize,
  linenos=true,
  frame=single,
  framesep=4pt,
  breaklines=true,
  breakanywhere=true,
  tabsize=4,
  numbersep=7pt,
  bgcolor=gray!8,
}
\setminted[python]{python3=true}

\SetupFloatingEnvironment{listing}{name=Лістинг,listname=Перелік лістингів}

% ---------------------------- Нумерація та оформлення заголовків ----------------------------
\usepackage{titlesec}

\setcounter{secnumdepth}{3}
\renewcommand{\thesection}{\arabic{section}}
\renewcommand{\thesubsection}{\arabic{section}.\arabic{subsection}}
\renewcommand{\thesubsubsection}{\arabic{section}.\arabic{subsection}.\arabic{subsubsection}}

% Розділ по центру, 16 пт, великими літерами; слово «РОЗДІЛ» опущене, бо звіт не дисертація. Номер лишається: \numberwithin прив'язує до нього нумерацію рисунків («Рис. 2.1»), і без нумерованих розділів вона б зламалася. \MakeUppercase четвертим аргументом отримує текст заголовка, тому в джерелі заголовки пишуться звичайним регістром.
\titleformat{\section}[block]
  {\centering\bfseries\fontsize{16}{19.2}\selectfont}
  {\thesection.}{0.3em}{\MakeUppercase}
% Підрозділи — зліва з абзацного відступу, звичайним регістром.
\titleformat{\subsection}[block]{\bfseries}{\thesubsection.}{1em}{}
\titleformat{\subsubsection}[block]{\bfseries\itshape}{\thesubsubsection.}{1em}{}

% Незіркові \titlespacing — обов'язкові. Без них titlesec застосовує зірковий варіант, який прибирає абзацний відступ у першого абзацу після заголовка й перекриває indentfirst. Значення — типові для report.
\titlespacing{\section}{0pt}{20pt}{6pt}
\titlespacing{\subsection}{1.27cm}{6pt}{0pt}
\titlespacing{\subsubsection}{1.27cm}{6pt}{0pt}

\setlength{\parindent}{1.27cm}
\setlength{\parskip}{0pt}
\renewcommand{\arraystretch}{1.3}

% Переноси слів вимкнені — так у шаблоні НУЛП. Натомість абзац розганяється міжслівними пробілами, тому tolerance і emergencystretch підняті: без них рядки виходили б за праве поле.
\hyphenpenalty=10000
\exhyphenpenalty=10000
\tolerance=9999
\setlength{\emergencystretch}{3em}

% ---------------------------- Підписи до рисунків і таблиць ----------------------------
\captionsetup[figure]{name=Рис.,justification=centering,font={bf,it}}
\captionsetup[table]{name=Таблиця,justification=raggedleft,singlelinecheck=false,labelfont=bf,textfont=normalfont}

% Рисунки, таблиці й формули нумеруються в межах розділу: «Рис. 2.1», а не «Рис. 5». Так у шаблоні НУЛП, і так само виглядали приклади на консультації.
\numberwithin{equation}{section}
\numberwithin{figure}{section}
\numberwithin{table}{section}

% ---------------------------- Колонтитули ----------------------------
\usepackage{fancyhdr}
\pagestyle{fancy}
\fancyhf{}
\fancyfoot[R]{\thepage}
\renewcommand{\headrulewidth}{0pt}
\fancypagestyle{plain}{%
  \fancyhf{}
  \fancyfoot[R]{\thepage}
  \renewcommand{\headrulewidth}{0pt}
}

\input{../../../../titlepage}
: \renewcommand{\WorkKind}{ПРАКТИЧНА РОБОТА}.
\newcommand{\WorkKind}{ЛАБОРАТОРНА РОБОТА}

\newcommand{\MakeLabTitle}[2]{%
\begin{titlepage}
% Титулка — не суцільний текст, а фіксована верстка з явних \\[...] проміжків,
% тож міжрядковий інтервал основного тексту (\setstretch{1.5} у preamble.tex)
% до неї не застосовний: він розтягував кожен із ~15 рядків на 20 %, і блок
% «Виконав» разом із «Львів --- рік» переповнював сторінку. Інтервал локальний
% для оточення titlepage, тому на решту документа це не впливає.
\setstretch{1}
\begin{center}
  {\large\bfseries МІНІСТЕРСТВО ОСВІТИ І НАУКИ УКРАЇНИ}\\[6pt]
  {\large\bfseries Національний університет «Львівська політехніка»}\\[6pt]
  {\large\bfseries Інститут комп'ютерних наук та інформаційних технологій}\\[6pt]
  {\large\bfseries Кафедра «Системи штучного інтелекту»}\\[1.2cm]

  \includegraphics[width=5cm]{../../../../logo}\\[1.2cm]

  {\Large\bfseries \WorkKind{} №#1}\\[0.5cm]
  з дисципліни\\[4pt]
  «\SubjectName»\\[0.6cm]
  на тему\\[4pt]
  \textit{#2}
\end{center}

% Відступ до блоку «Виконав» гнучкий, а не фіксований (було 2 cm). Фіксований
% разом із довгим родом роботи чи темою на два рядки переповнював сторінку, і
% «Львів --- рік» зі свого \vfill з'їжджав на окрему другу сторінку. \vfill
% стискається до нуля, коли місця бракує, тож переповнення стало неможливим.
\vfill

\begin{flushright}
  {\bfseries Виконав:}\\
  студент групи \StudentGroup\\
  \StudentName\\[1cm]
  {\bfseries Викладач:}\\
  \LecturerName\\
\end{flushright}

\vfill
\begin{center}
  Львів --- \AcademicYear
\end{center}
\end{titlepage}
\clearpage
}


% Subject metadata, read by the shared title page (titlepage.tex).

\newcommand{\SubjectName}{Методологія та процеси моделювання у штучному інтелекті}
% The lecturer's position and degree are not known yet; the other subjects put them on a line above the name, separated by `\\`.
\newcommand{\LecturerName}{Косаревич Р.Я.}
\newcommand{\StudentName}{Галич Максим}
\newcommand{\StudentGroup}{КНСШ-13}
\newcommand{\AcademicYear}{2026}

% Figures are generated artifacts, so they live in .cache/ rather than beside the text (see ARTIFACTS in src/mk1/config.py).
\graphicspath{{../../../../.cache/reports-artifacts/1th term 2026 Autumn/Методологія та процеси моделювання у штучному інтелекті/}}

\begin{document}

\MakeLabTitle{2}{Порівняння традиційних ML-алгоритмів для задачі класифікації}

\textbf{Мета роботи.} Ознайомитися з основними алгоритмами машинного навчання та навчитися обирати й порівнювати їхню ефективність для конкретної задачі.

Усі обчислення виконано мовою Python із бібліотеками scikit-learn та imbalanced-learn. Код продовжує пакет \texttt{mk1} з лабораторної роботи №1: сценарії підготовки даних і моделі описано в модулі \texttt{scenarios.py}, порівняння --- у \texttt{comparison.py}, дослідження перенавчання --- в \texttt{overfitting.py}. Сценарій \texttt{lab2/run.py} лише запускає етапи по черзі й перевіряє, що кожен записав свої рисунки та метрики. Увесь експеримент відтворюється однією командою з фіксованим зерном генератора випадкових чисел.

\section{Постановка задачі та вихідні дані}

Використано набір даних, підготовлений у лабораторній роботі №1: \textbf{Diabetes 130-US Hospitals for Years 1999--2008} [\ref{ref:strack}]. Після очищення він містить \textbf{69\,987} незалежних пацієнтів (для кожного залишено першу госпіталізацію) і 46 числових та категоріальних ознак. Цільова змінна --- повторна госпіталізація протягом 30 днів після виписки, тож задача є \textbf{бінарною класифікацією}. Клас 1 становить 8,98~\% спостережень, співвідношення класів 1~:~10,1.

Розбиття на навчальну та тестову вибірки взято без змін з лабораторної роботи №1: стратифіковане, 80~\% на 20~\%, з тим самим зерном. Навчальна вибірка містить 55\,989 пацієнтів, тестова --- 13\,998, з яких 1\,257 були повторно госпіталізовані. Усі порівняння в розділах~\ref{sec:results}--\ref{sec:overfitting} виконано крос-валідацією на навчальній вибірці. Тестову вибірку використано один раз, у розділі~\ref{sec:test}, для перевірки висновків.

\section{Сценарії підготовки даних}
\label{sec:scenarios}

\subsection{Опис сценаріїв}

Порівняння виконано на трьох сценаріях підготовки даних, які вимагає завдання (табл.~\ref{tab:scenarios}). Сценарій В виконано у двох варіантах: з методом SMOTE [\ref{ref:smote}] і з методом ADASYN [\ref{ref:adasyn}].

% A longtable rather than a table float: rows of prose may run longer than the rest of a page, and an [H] float taller than a page sends LaTeX into an endless loop of empty pages.
{\small
\renewcommand{\arraystretch}{1.15}
\begin{longtable}{|p{3.0cm}|p{4.3cm}|p{4.3cm}|p{3.2cm}|}
\caption{Перетворення даних у сценаріях порівняння}
\label{tab:scenarios} \\
\hline
\textbf{Етап} & \textbf{А: мінімальне очищення} & \textbf{Б: повне очищення} & \textbf{В: очищення + аугментація} \\
\hline
\endfirsthead
\hline
\textbf{Етап} & \textbf{А: мінімальне очищення} & \textbf{Б: повне очищення} & \textbf{В: очищення + аугментація} \\
\hline
\endhead
Записи & \multicolumn{3}{p{12.2cm}|}{однакові в усіх сценаріях: перша госпіталізація кожного пацієнта, без випадків смерті та переведення до хоспісу; службові коди «NULL» і «Not Mapped» замінено на пропуски} \\
\hline
Пропуски категоріальних ознак & найчастіше значення (мода) & окрема категорія «Missing» & як у Б \\
\hline
Високопотужні категоріальні ознаки & One-Hot з об'єднанням категорій, рідших за 50 записів & Target Encoding з перехресним навчанням [\ref{ref:target}] & як у Б \\
\hline
Інші категоріальні ознаки & One-Hot & One-Hot & як у Б \\
\hline
Вік & порядкове кодування & порядкове кодування, стандартизація & як у Б \\
\hline
Числові ознаки & без перетворень & \texttt{RobustScaler} & як у Б \\
\hline
Інженерія ознак & немає & доменні ознаки, групи діагнозів МКХ-9, відхилення від середнього за категорією & як у Б \\
\hline
Балансування класів & немає & немає & SMOTE або ADASYN до співвідношення 1~:~1, лише на навчальних блоках \\
\hline
\end{longtable}
}

Записи в усіх сценаріях однакові. Очищення на рівні записів не є вибором методу обробки, а умовою коректного оцінювання. Без нього госпіталізації одного пацієнта потрапили б і до навчальної, і до валідаційної частини, і метрики сценарію А були б завищені витоком даних, а не відображали б якість обробки. Отже, сценарії відрізняються лише тим, що зроблено з ознаками, і різниця в метриках є наслідком саме цього.

\textbf{Сценарій А} --- це мінімум обробки, без якого бібліотека scikit-learn не прийме дані: пропуски заповнено модою, категорії закодовано One-Hot, числові ознаки подано як є (лістинг~\ref{lst:minimal}). Рідкісні категорії все ж об'єднано: без цього лише коди діагнозів дали б близько 2000 стовпців, переважно порожніх.

\inputminted[firstline=92,lastline=111]{python}{../src/mk1/scenarios.py}
\captionof{listing}{Мінімальна обробка сценарію А (\texttt{src/mk1/scenarios.py})}
\label{lst:minimal}

\textbf{Сценарій Б} --- підсумковий конвеєр лабораторної роботи №1, кожен крок якого обрано за результатами крос-валідації. Пропуски категоріальних ознак заповнено окремою категорією «Missing», а не імп'ютацією: у лабораторній роботі №1 цей спосіб переміг моду та імп'ютацію методом $k$ найближчих сусідів, бо сам факт пропуску пов'язаний із ціллю (пропуск коду платника чи спеціальності лікаря --- ознака певного типу госпіталізації). Числові ознаки цього набору даних пропусків не мають. Для всіх трьох моделей використано одну конфігурацію, щоб сценарій був властивістю даних, а не моделі. У лабораторній роботі №1 деревоподібні моделі були байдужі до вибору кодування та масштабування в межах розкиду між блоками.

\textbf{Сценарій В} додає до сценарію Б передискретизацію класу 1. SMOTE створює нову точку на відрізку між записом класу 1 та одним із його п'яти найближчих сусідів того ж класу. ADASYN генерує більше точок біля записів, серед сусідів яких більше записів класу 0, тобто біля межі класів. Обидва методи працюють у закодованому числовому просторі, після кроку попередньої обробки.

Кожен сценарій зібрано в один конвеєр бібліотеки imbalanced-learn (лістинг~\ref{lst:build}). Усі кроки, що навчаються на даних (групові середні, кодувальники, масштабування, передискретизація), стоять усередині конвеєра. Під час крос-валідації їх заново навчають на навчальних блоках, а до валідаційного блоку лише застосовують. Передискретизація виконується тільки під час навчання, тож оцінювання завжди відбувається на реальних пацієнтах.

\inputminted[firstline=128,lastline=160]{python}{../src/mk1/scenarios.py}
\captionof{listing}{Класифікатори та збирання конвеєра сценарію (\texttt{src/mk1/scenarios.py})}
\label{lst:build}

\subsection{Стан даних після перетворень}

Стан навчальної матриці, яку отримує класифікатор у кожному сценарії, наведено в табл.~\ref{tab:profiles} (для всієї навчальної вибірки).

\begin{table}[H]
\caption{Навчальна матриця після перетворень кожного сценарію}
\label{tab:profiles}
\begin{tabularx}{\linewidth}{|X|r|r|r|r|}
\hline
\textbf{Сценарій} & \textbf{Записів} & \textbf{З них синтетичних} & \textbf{Частка класу 1} & \textbf{Ознак} \\
\hline
А: мінімальне очищення & 55\,989 & 0 & 8,98~\% & 563 \\
\hline
Б: повне очищення & 55\,989 & 0 & 8,98~\% & 113 \\
\hline
В: очищення + SMOTE & 101\,922 & 45\,933 & 50,00~\% & 113 \\
\hline
В: очищення + ADASYN & 100\,307 & 44\,318 & 49,19~\% & 113 \\
\hline
\end{tabularx}
\end{table}

Сценарій А дає найширшу матрицю: One-Hot кодування високопотужних ознак перетворює сім стовпців на сотні бінарних. Target Encoding у сценарії Б замінює кожну з цих ознак одним числом, тож матриця вужча в п'ять разів, попри додані доменні ознаки. Аугментація в сценарії В майже подвоює кількість записів: на кожного реального пацієнта класу 1 припадає близько дев'яти синтетичних.

\section{Алгоритми класифікації}
\label{sec:models}

Для порівняння обрано три алгоритми з різними принципами роботи (лістинг~\ref{lst:build}, функція \texttt{make\_classifier}):

\begin{itemize}
\item \textbf{логістична регресія} (L2-регуляризація, $C = 1$, розв'язувач L-BFGS, до 3000 ітерацій): лінійна модель, що оцінює ймовірність класу 1 логістичною функцією від зваженої суми ознак;
\item \textbf{випадковий ліс} [\ref{ref:breiman}] (300 дерев, $\sqrt{p}$ випадкових ознак у кожному вузлі, щонайменше 10 записів у листку): ансамбль дерев рішень, навчених на бутстреп-вибірках; прогноз --- середнє прогнозів дерев;
\item \textbf{гістограмний градієнтний бустинг} [\ref{ref:friedman}, \ref{ref:lightgbm}] (250 дерев, крок навчання 0,05, до 31 листка в дереві, L2-регуляризація 1): дерева будуються послідовно, кожне виправляє похибки попередніх; значення ознак попередньо групуються в 255 інтервалів, що прискорює пошук розбиттів.
\end{itemize}

Гіперпараметри логістичної регресії та бустингу взято з лабораторної роботи №1. Мінімальний розмір листка випадкового лісу обрано за найвищою PR-AUC на крос-валідації серед значень 1, 3, 10, 30 і 100 (розділ~\ref{sec:capacity}).

Метод опорних векторів (SVM) до порівняння не включено. Час навчання SVM з ядром зростає приблизно пропорційно квадрату кількості записів або швидше [\ref{ref:libsvm}], а в цьому експерименті модель довелося б навчити понад сто разів на 45--100 тисячах записів. Крім того, SVM не дає ймовірностей без додаткової калібровки на окремій крос-валідації, а саме ймовірності потрібні для ROC-AUC, PR-AUC та вибору порога.

\section{Методика оцінювання}
\label{sec:method}

Кожну пару «сценарій --- модель» оцінено стратифікованою 5-блоковою крос-валідацією на навчальній вибірці (лістинг~\ref{lst:cv}) і подано як середнє $\pm$ стандартне відхилення між блоками.

\inputminted[firstline=119,lastline=149]{python}{../src/mk1/comparison.py}
\captionof{listing}{Крос-валідація з оцінкою на навчальних даних (\texttt{src/mk1/comparison.py})}
\label{lst:cv}

Метрики класифікації, яких вимагає завдання, обчислено для класу 1: accuracy (частка правильних відповідей), precision (частка справжніх випадків серед позначених моделлю), recall (частка знайдених випадків серед усіх справжніх), F1-міра та ROC-AUC. Додатково наведено PR-AUC --- площу під кривою точність-повнота, основну метрику лабораторної роботи №1, що краще за ROC-AUC відображає якість на незбалансованих даних [\ref{ref:saito}], та оцінку Брієра --- середній квадрат різниці між імовірністю та фактичним класом, що вимірює якість самих імовірностей.

Accuracy, precision, recall і F1 залежать від порога рішення, тому їх обчислено за двох порогів:

\begin{itemize}
\item \textbf{стандартний поріг 0,5}, який використовує метод \texttt{predict} бібліотеки scikit-learn;
\item \textbf{оптимальний поріг}, за якого F1-міра на прогнозах крос-валідації (out-of-fold) максимальна. Для тестової вибірки поріг береться саме звідси, тобто обирається лише на навчальних даних.
\end{itemize}

Перенавчання виміряно розривом між якістю на навчальних даних і на валідаційних блоках. Якість на навчальних даних обчислено лише на реальних записах навчальної матриці, без синтетичних. Для сценарію В додатково обчислено частку синтетичних записів, які модель відносить до класу 1 за порога 0,5. Якщо модель навчилася розпізнавати синтетичні точки, а не пацієнтів, ця частка буде близькою до 1, а повнота на реальних пацієнтах --- значно нижчою. Ті самі навчальні записи також оцінено так, як оцінюється новий пацієнт, тобто пропущено через навчений конвеєр. Різниця між двома оцінками виникає через Target Encoding: рядки навчальної матриці закодовано з перехресним навчанням, а новий пацієнт отримує кодування, навчене на всій навчальній частині.

\section{Результати порівняння}
\label{sec:results}

Значення всіх метрик для кожного сценарію та моделі наведено на рис.~\ref{fig:comparison}, а числові результати --- у табл.~\ref{tab:cv_05}--\ref{tab:cv_tuned}.

\begin{figure}[H]
  \centering
  \includegraphics[width=\linewidth]{lab2/comparison_metrics}
  \caption{Метрики на валідаційних блоках крос-валідації (середнє $\pm$ стандартне відхилення за 5 блоками)}
  \label{fig:comparison}
\end{figure}

\subsection{Метрики за стандартного порога}

\begin{table}[H]
\caption{Крос-валідація: рішення за порога 0,5}
\label{tab:cv_05}
\small
\renewcommand{\arraystretch}{1.15}
\begin{tabularx}{\linewidth}{|X|c|c|c|c|}
\hline
\textbf{Модель} & \textbf{Accuracy} & \textbf{Precision} & \textbf{Recall} & \textbf{F1} \\
\hline
\multicolumn{5}{|l|}{\textit{А: мінімальне очищення}} \\
\hline
логістична регресія & 0,910 & 0,430 & 0,008 & 0,016 \\
\hline
випадковий ліс & 0,910 & 0,000 & 0,000 & 0,000 \\
\hline
градієнтний бустинг & 0,910 & 0,547 & 0,006 & 0,012 \\
\hline
\multicolumn{5}{|l|}{\textit{Б: повне очищення}} \\
\hline
логістична регресія & 0,910 & 0,431 & 0,007 & 0,014 \\
\hline
випадковий ліс & 0,910 & 0,000 & 0,000 & 0,000 \\
\hline
градієнтний бустинг & 0,910 & 0,557 & 0,006 & 0,012 \\
\hline
\multicolumn{5}{|l|}{\textit{В: очищення + SMOTE}} \\
\hline
логістична регресія & 0,633 & 0,131 & 0,550 & 0,212 \\
\hline
випадковий ліс & 0,905 & 0,196 & 0,015 & 0,026 \\
\hline
градієнтний бустинг & 0,758 & 0,116 & 0,231 & 0,145 \\
\hline
\multicolumn{5}{|l|}{\textit{В: очищення + ADASYN}} \\
\hline
логістична регресія & 0,635 & 0,131 & 0,544 & 0,211 \\
\hline
випадковий ліс & 0,906 & 0,220 & 0,012 & 0,022 \\
\hline
градієнтний бустинг & 0,797 & 0,158 & 0,187 & 0,133 \\
\hline
\end{tabularx}
\end{table}

У сценаріях А і Б accuracy усіх трьох моделей дорівнює 0,910, тобто рівно частці класу 0. Моделі майже нікого не відносять до класу 1: recall не перевищує 0,008, а випадковий ліс не позначив жодного пацієнта. Модель, навчена на вибірці з 9~\% класу 1, рідко видає ймовірність понад 0,5. Для випадкового лісу це ще помітніше, бо його прогноз --- середнє голосів 300 дерев, і ймовірності зсунуті від крайніх значень до середини [\ref{ref:niculescu}]. За порога 0,5 рішення всіх моделей майже збігаються з тривіальним правилом «завжди клас 0», яке теж має accuracy 0,910. Тому accuracy у цій задачі непридатна для порівняння моделей, а F1-міра за порога 0,5 показує не якість моделі, а невідповідність порога розподілу її ймовірностей.

Аугментація змінює картину. SMOTE піднімає recall логістичної регресії до 0,550, а F1 --- з 0,014 до 0,212, хоча accuracy падає до 0,633. На перший погляд аугментація суттєво покращила модель. Насправді вона лише зсунула ймовірності вгору: модель, навчена на вибірці 1~:~1, вважає повторну госпіталізацію подією з шансом близько 50~\%, а не 9~\%. Той самий ефект можна отримати без аугментації, просто обравши інший поріг [\ref{ref:elor}] (підрозділ~\ref{sec:tuned}). Для випадкового лісу аугментація майже нічого не змінила за порога 0,5 (recall 0,015), хоча модель навчалася на вибірці 1~:~1. Причину цього розкрито в розділі~\ref{sec:overfitting}.

\subsection{Якість ранжування}

Метрики, що не залежать від порога, наведено в табл.~\ref{tab:cv_rank}. Вони показують, наскільки добре модель упорядковує пацієнтів за ризиком, і саме за ними коректно порівнювати сценарії.

\begin{table}[H]
\caption{Крос-валідація: якість ранжування та ймовірностей}
\label{tab:cv_rank}
\small
\renewcommand{\arraystretch}{1.15}
\begin{tabularx}{\linewidth}{|X|c|c|c|}
\hline
\textbf{Модель} & \textbf{ROC-AUC} & \textbf{PR-AUC} & \textbf{Брієр} \\
\hline
\multicolumn{4}{|l|}{\textit{А: мінімальне очищення}} \\
\hline
логістична регресія & 0,641 ± 0,005 & 0,166 ± 0,003 & 0,080 \\
\hline
випадковий ліс & 0,658 ± 0,007 & 0,178 ± 0,004 & 0,079 \\
\hline
градієнтний бустинг & 0,652 ± 0,008 & 0,180 ± 0,007 & 0,079 \\
\hline
\multicolumn{4}{|l|}{\textit{Б: повне очищення}} \\
\hline
логістична регресія & 0,651 ± 0,006 & 0,172 ± 0,005 & 0,079 \\
\hline
випадковий ліс & 0,653 ± 0,008 & 0,177 ± 0,004 & 0,079 \\
\hline
градієнтний бустинг & 0,655 ± 0,006 & 0,179 ± 0,005 & 0,079 \\
\hline
\multicolumn{4}{|l|}{\textit{В: очищення + SMOTE}} \\
\hline
логістична регресія & 0,636 ± 0,007 & 0,163 ± 0,006 & 0,228 \\
\hline
випадковий ліс & 0,622 ± 0,021 & 0,136 ± 0,009 & 0,111 \\
\hline
градієнтний бустинг & 0,537 ± 0,020 & 0,112 ± 0,009 & 0,175 \\
\hline
\multicolumn{4}{|l|}{\textit{В: очищення + ADASYN}} \\
\hline
логістична регресія & 0,634 ± 0,006 & 0,161 ± 0,006 & 0,226 \\
\hline
випадковий ліс & 0,617 ± 0,022 & 0,131 ± 0,010 & 0,112 \\
\hline
градієнтний бустинг & 0,560 ± 0,049 & 0,123 ± 0,021 & 0,152 \\
\hline
\end{tabularx}
\end{table}

\textbf{Повне очищення (А $\to$ Б) покращило лише логістичну регресію}: ROC-AUC зросла з 0,641 до 0,651, PR-AUC --- з 0,166 до 0,172, що перевищує розкид між блоками. Для деревоподібних моделей різниця між сценаріями менша за одне стандартне відхилення: випадковий ліс має PR-AUC 0,178 і 0,177, бустинг --- 0,180 і 0,179. Лінійна модель залежить від кожного кроку обробки: без масштабування штраф регуляризації нерівномірно діє на ознаки різного масштабу, а кожна з сотень One-Hot ознак рідкісних категорій отримує власний коефіцієнт. Дерева інваріантні до монотонних перетворень ознак і самі будують розбиття за окремими категоріями, тож масштабування й кодування для них майже нічого не змінюють.

Повне очищення також у 21 раз прискорило навчання логістичної регресії: 1,4~с проти 29,4~с на блок. Немасштабовані лічильники (до 132 аналізів за госпіталізацію) роблять задачу оптимізації погано обумовленою, і розв'язувачу L-BFGS потрібно значно більше ітерацій.

\textbf{Аугментація погіршила всі моделі.} Для логістичної регресії втрата невелика: PR-AUC 0,163 і 0,161 проти 0,172. Для деревоподібних моделей вона суттєва: випадковий ліс втратив 0,041--0,046 PR-AUC, бустинг --- 0,056--0,067, а ROC-AUC бустингу з SMOTE (0,537) майже дорівнює рівню випадкового вгадування. Розкид між блоками для деревоподібних моделей з аугментацією зростає у 2--8 разів, тобто результат стає ще й нестабільним.

Оцінка Брієра без аугментації становить 0,079--0,080, що нижче, ніж у сталого прогнозу базової частоти ($0{,}090 \cdot 0{,}910 = 0{,}082$): імовірності моделей відкалібровані. З аугментацією вона зростає до 0,111--0,228, тобто ймовірності перестають відповідати фактичній частоті повторних госпіталізацій. Для моделі клінічного ризику це самостійна втрата [\ref{ref:goorbergh}].

Найвищу якість ранжування на крос-валідації мають деревоподібні моделі: бустинг (PR-AUC 0,180 у сценарії А і 0,179 у Б) та випадковий ліс (0,178 і 0,177). Різниця між ними менша за розкид між блоками. Логістична регресія відстає на 0,005--0,008 PR-AUC навіть після повного очищення.

\subsection{Метрики за оптимального порога}
\label{sec:tuned}

\begin{table}[H]
\caption{Крос-валідація: рішення за порога, що максимізує F1-міру}
\label{tab:cv_tuned}
\small
\renewcommand{\arraystretch}{1.15}
\begin{tabularx}{\linewidth}{|X|c|c|c|c|}
\hline
\textbf{Модель} & \textbf{Поріг} & \textbf{Precision} & \textbf{Recall} & \textbf{F1} \\
\hline
\multicolumn{5}{|l|}{\textit{А: мінімальне очищення}} \\
\hline
логістична регресія & 0,121 & 0,163 & 0,359 & 0,224 \\
\hline
випадковий ліс & 0,120 & 0,175 & 0,354 & 0,234 \\
\hline
градієнтний бустинг & 0,105 & 0,154 & 0,460 & 0,231 \\
\hline
\multicolumn{5}{|l|}{\textit{Б: повне очищення}} \\
\hline
логістична регресія & 0,103 & 0,156 & 0,432 & 0,229 \\
\hline
випадковий ліс & 0,117 & 0,168 & 0,400 & 0,237 \\
\hline
градієнтний бустинг & 0,134 & 0,193 & 0,296 & 0,233 \\
\hline
\multicolumn{5}{|l|}{\textit{В: очищення + SMOTE}} \\
\hline
логістична регресія & 0,587 & 0,159 & 0,345 & 0,218 \\
\hline
випадковий ліс & 0,296 & 0,140 & 0,440 & 0,204 \\
\hline
градієнтний бустинг & 0,140 & 0,092 & 0,943 & 0,168 \\
\hline
\multicolumn{5}{|l|}{\textit{В: очищення + ADASYN}} \\
\hline
логістична регресія & 0,601 & 0,169 & 0,300 & 0,216 \\
\hline
випадковий ліс & 0,263 & 0,125 & 0,594 & 0,204 \\
\hline
градієнтний бустинг & 0,257 & 0,111 & 0,624 & 0,181 \\
\hline
\end{tabularx}
\end{table}

Без аугментації оптимальний поріг лежить у межах 0,10--0,13, близько до базової частоти класу 1, і дає F1-міру 0,224--0,237. З аугментацією поріг логістичної регресії зсувається до 0,59--0,60, а F1-міра становить 0,218 і 0,216, що нижче, ніж у сценарії Б (0,229). Отже, для лінійної моделі вся видима користь аугментації за порога 0,5 зводиться до зсуву порога, а за правильно обраного порога аугментація дає невеликий програш. Для деревоподібних моделей F1-міра з аугментацією падає до 0,17--0,20. Найвищу F1-міру має випадковий ліс у сценарії Б: 0,237 за precision 0,168 і recall 0,400.

\section{Аналіз перенавчання та впливу аугментації}
\label{sec:overfitting}

\subsection{Розрив між навчальними та валідаційними даними}

Якість на навчальних даних і на валідаційних блоках порівняно на рис.~\ref{fig:gap} і в табл.~\ref{tab:gap}.

\begin{figure}[H]
  \centering
  \includegraphics[width=\linewidth]{lab2/overfitting_gap}
  \caption{Якість на навчальних даних (квадрати) і на валідаційних блоках (кола)}
  \label{fig:gap}
\end{figure}

\begin{table}[H]
\caption{Перенавчання: ROC-AUC та частка класу 1, розпізнана за порога 0,5}
\label{tab:gap}
\footnotesize
\renewcommand{\arraystretch}{1.15}
\setlength{\tabcolsep}{4pt}
\begin{tabularx}{\linewidth}{|X|c|c|c|c|c|c|c|}
\hline
 & \multicolumn{3}{c|}{\textbf{ROC-AUC}} & \multicolumn{4}{c|}{\textbf{Recall за порога 0,5}} \\
\cline{2-8}
\textbf{Модель} & \textbf{навч.} & \textbf{вал.} & \textbf{розрив} & \textbf{навч.} & \textbf{навч. як нові} & \textbf{вал.} & \textbf{синтет.} \\
\hline
\multicolumn{8}{|l|}{\textit{А: мінімальне очищення}} \\
\hline
логістична регресія & 0,690 & 0,641 & 0,049 & 0,009 & 0,009 & 0,008 & --- \\
\hline
випадковий ліс & 0,878 & 0,658 & 0,220 & 0,000 & 0,000 & 0,000 & --- \\
\hline
градієнтний бустинг & 0,816 & 0,652 & 0,163 & 0,020 & 0,020 & 0,006 & --- \\
\hline
\multicolumn{8}{|l|}{\textit{Б: повне очищення}} \\
\hline
логістична регресія & 0,661 & 0,651 & 0,010 & 0,007 & 0,010 & 0,007 & --- \\
\hline
випадковий ліс & 0,985 & 0,653 & 0,332 & 0,000 & 0,000 & 0,000 & --- \\
\hline
градієнтний бустинг & 0,870 & 0,655 & 0,215 & 0,035 & 0,026 & 0,006 & --- \\
\hline
\multicolumn{8}{|l|}{\textit{В: очищення + SMOTE}} \\
\hline
логістична регресія & 0,648 & 0,636 & 0,012 & 0,568 & 0,610 & 0,550 & 0,614 \\
\hline
випадковий ліс & 0,986 & 0,622 & 0,364 & 0,102 & 0,182 & 0,015 & 1,000 \\
\hline
градієнтний бустинг & 0,809 & 0,537 & 0,272 & 0,018 & 0,369 & 0,231 & 0,999 \\
\hline
\multicolumn{8}{|l|}{\textit{В: очищення + ADASYN}} \\
\hline
логістична регресія & 0,646 & 0,634 & 0,012 & 0,563 & 0,602 & 0,544 & 0,589 \\
\hline
випадковий ліс & 0,988 & 0,617 & 0,371 & 0,081 & 0,165 & 0,012 & 1,000 \\
\hline
градієнтний бустинг & 0,810 & 0,560 & 0,250 & 0,018 & 0,295 & 0,187 & 0,999 \\
\hline
\end{tabularx}
\end{table}

\textbf{Логістична регресія} майже не перенавчається: розрив ROC-AUC у сценаріях Б і В становить 0,010--0,012. У сценарії А він уп'ятеро більший (0,049): кожна з майже 500 One-Hot ознак, зокрема рідкісні коди діагнозів, має власний коефіцієнт, і модель частково запам'ятовує шум. Target Encoding стискає кожну високопотужну ознаку до одного числа, тож для лінійної моделі правильне кодування діє як регуляризація: розрив зменшується, а якість на валідації зростає.

\textbf{Деревоподібні моделі} мають великий розрив навіть без аугментації: випадковий ліс 0,220 і 0,332, бустинг 0,163 і 0,215. Проте саме вони мають найкращу якість на валідації. Розрив вимірює, наскільки модель запам'ятала навчальні записи, але не те, чи шкодить їй це. Кожне дерево лісу запам'ятовує свою бутстреп-вибірку, а усереднення 300 дерев гасить більшу частину цього шуму [\ref{ref:breiman}].

З аугментацією розрив деревоподібних моделей зростає (випадковий ліс 0,364 і 0,371, бустинг 0,272 і 0,250), а якість на валідації падає. Ключ до пояснення дають останні стовпці табл.~\ref{tab:gap}. Випадковий ліс відносить до класу 1 100~\% синтетичних записів, бустинг --- 99,9~\%. Водночас серед реальних повторно госпіталізованих пацієнтів, на яких ці моделі навчалися, вони розпізнають лише 10,2~\% і 1,8~\% відповідно. \textbf{Моделі навчилися відрізняти синтетичні записи від реальних, а не пацієнтів із ризиком від решти.} Логістична регресія так не може: вона відносить до класу 1 61~\% синтетичних записів, 57~\% реальних навчальних і 55~\% валідаційних, тобто сприймає синтетичні записи так само, як реальні.

\subsection{Чому деревоподібні моделі розпізнають синтетичні записи}

Синтетичний запис SMOTE лежить на відрізку між двома реальними записами в закодованому просторі ознак. У цьому наборі даних такий запис має значення, яких не буває в реальних пацієнтів:

\begin{itemize}
\item \textbf{дробові значення One-Hot ознак}. Якщо два сусіди мають різні категорії, синтетичний запис отримує, наприклад, 0,37 у стовпці категорії, а реальні пацієнти мають там лише 0 або 1. Дереву досить двох розбиттів, щоб виділити всі такі записи. Змагальна перевірка в лабораторній роботі №1 відрізняла синтетичні записи від реальних з ROC-AUC 0,98--1,00;
\item \textbf{проміжні значення Target Encoding}. Через перехресне навчання кожна категорія в навчальній матриці має лише кілька різних значень, по одному на блок. Синтетичні записи заповнюють проміжки між ними. Дерево вивчає правило «точне значення реального запису --- клас 0, проміжне значення --- клас 1».
\end{itemize}

Друге правило руйнує прогнози для нових пацієнтів. Новий пацієнт отримує кодування, навчене на всій навчальній частині, і воно не збігається з жодним значенням навчальної матриці. Тому для дерева новий пацієнт схожий на синтетичний запис. Це підтверджує стовпець «навч. як нові» табл.~\ref{tab:gap}: ті самі реальні навчальні пацієнти, закодовані так, як кодується новий пацієнт, відносяться бустингом до класу 1 у 36,9~\% випадків замість 1,8~\%, а випадковим лісом --- у 18,2~\% замість 10,2~\%. Бустинг фактично ділить пацієнтів на «схожих на навчальні» та «схожих на синтетичні», і цей поділ майже не пов'язаний із ризиком повторної госпіталізації. Звідси ROC-AUC 0,537. У сценаріях без аугментації синтетичних записів немає, і два способи кодування дають майже однакові прогнози 2,6~\% проти 3,5~\% для бустингу.

\subsection{Складність моделі}
\label{sec:capacity}

Щоб перевірити, чи шкода від аугментації є звичайним перенавчанням надто складної моделі, мінімальний розмір листка випадкового лісу змінено від 1 (дерева ростуть до повної чистоти листків) до 100 (неглибокі дерева) у сценаріях Б і В (рис.~\ref{fig:capacity}, лістинг~\ref{lst:sweeps}).

\inputminted[firstline=27,lastline=45]{python}{../src/mk1/overfitting.py}
\captionof{listing}{Параметри досліджень складності моделі та обсягу аугментації (\texttt{src/mk1/overfitting.py})}
\label{lst:sweeps}

\begin{figure}[H]
  \centering
  \includegraphics[width=\linewidth]{lab2/overfitting_capacity}
  \caption{Випадковий ліс: якість на навчальних даних і на валідації залежно від мінімального розміру листка}
  \label{fig:capacity}
\end{figure}

Без аугментації залежність класична. Дерева з листками по одному запису ідеально запам'ятовують навчальні дані (ROC-AUC 1,000), але мають найгіршу якість на валідації: PR-AUC 0,166, ROC-AUC 0,636. Зі збільшенням листка навчальна ROC-AUC падає до 0,759, а валідаційна зростає до 0,658. PR-AUC максимальна за листка 10 (0,177), тому це значення використано в основному порівнянні.

З аугментацією SMOTE валідаційна PR-AUC лишається в межах 0,132--0,137 за будь-якого розміру листка. Навіть неглибокі дерева з листками по 100 записів відносять до класу 1 97,6~\% синтетичних записів. Отже, \textbf{регуляризація моделі не усуває шкоди від аугментації}: модель вивчає не випадковий шум окремих записів, а систематичну відмінність синтетичних записів від реальних, і для цього їй не потрібні глибокі дерева.

\subsection{Обсяг аугментації}

Другий параметр, від якого залежить шкода, --- кількість синтетичних записів. SMOTE застосовано з цільовим співвідношенням класів від 0,25 до 1,0. Початкова точка --- власне співвідношення даних 0,099 без аугментації (рис.~\ref{fig:ratio}).

\begin{figure}[H]
  \centering
  \includegraphics[width=\linewidth]{lab2/overfitting_ratio}
  \caption{Вплив обсягу аугментації SMOTE на якість, розрив і розпізнавання синтетичних записів}
  \label{fig:ratio}
\end{figure}

Якість обох деревоподібних моделей монотонно погіршується зі збільшенням частки синтетичних записів. PR-AUC випадкового лісу падає з 0,177 до 0,165, 0,150, 0,141 і 0,136, а розрив ROC-AUC зростає з 0,332 до 0,364. Бустинг втрачає більшу частину якості вже за найменшого співвідношення 0,25: PR-AUC падає з 0,179 до 0,135, і вже тоді він відносить до класу 1 98,9~\% синтетичних записів. Для логістичної регресії втрата PR-AUC невелика (з 0,172 до 0,163), розрив не змінюється (0,010--0,012), а частка синтетичних записів, віднесених до класу 1, щоразу близька до recall на реальних валідаційних пацієнтах (від 0,051 і 0,055 до 0,614 і 0,550). Для неї синтетичні записи не відрізняються від реальних, і аугментація лише зсуває ймовірності: оцінка Брієра зростає з 0,079 до 0,228.

\subsection{Підсумок: чи покращила аугментація якість передбачення}

\textbf{Аугментація не покращила жодної моделі} за жодною метрикою, що не залежить від порога, ні на крос-валідації, ні на тестовій вибірці (розділ~\ref{sec:test}). Єдине видиме покращення --- зростання recall і F1-міри логістичної регресії за порога 0,5 --- є наслідком зсуву ймовірностей. Обравши поріг без аугментації, можна отримати вищу F1-міру (0,229 проти 0,218) і відкалібровані ймовірності. Це узгоджується з висновками інших досліджень про те, що за сильних класифікаторів і правильно обраного порога балансування класів не дає переваги [\ref{ref:elor}, \ref{ref:goorbergh}].

Проведені експерименти дозволяють сформулювати умови, за яких аугментація призводить до перенавчання:

\begin{enumerate}
\item \textbf{Гнучка модель.} Деревоподібні моделі здатні виділити невелику область простору ознак, де лежать синтетичні записи. Лінійна модель цього не може і тому від аугментації майже не страждає.
\item \textbf{Синтетичні записи відрізняються від реальних.} Інтерполяція між дискретними значеннями (лічильники, One-Hot категорії, закодовані значення) дає неможливі для реальних пацієнтів значення. На неперервних ознаках із плавною межею класів SMOTE безпечніший.
\item \textbf{Велика частка синтетичних записів.} За співвідношення 1~:~1 синтетичних записів класу 1 приблизно в дев'ять разів більше, ніж реальних, і клас 1 для моделі фактично стає класом синтетичних записів. Шкода зростає з кожним кроком обсягу аугментації.
\item \textbf{Різне кодування навчальних і нових записів.} Кодування з перехресним навчанням перед аугментацією робить нових пацієнтів схожими на синтетичних.
\item \textbf{Аугментація до розбиття на навчальну та валідаційну частини.} У цій роботі передискретизація стоїть усередині конвеєра й бачить лише навчальні блоки. Якби SMOTE застосували до всієї вибірки перед крос-валідацією, синтетичні сусіди валідаційних записів потрапили б до навчальних даних, і оцінка якості була б завищеною, а не заниженою, як тут.
\end{enumerate}

Важливий методологічний висновок: сам розрив між навчальною та валідаційною якістю не виявляє цієї проблеми. Для випадкового лісу він зріс лише з 0,332 до 0,364, і без аугментації був великим теж. Проблему виявляє порівняння частки розпізнаних синтетичних записів із recall на реальних пацієнтах.

\section{Сильні та слабкі сторони алгоритмів}
\label{sec:strengths}

Сильні та слабкі сторони кожного алгоритму, виявлені в експериментах, зведено в табл.~\ref{tab:strengths}.

{\small
\renewcommand{\arraystretch}{1.15}
\begin{longtable}{|p{2.8cm}|p{6.4cm}|p{6.4cm}|}
\caption{Сильні та слабкі сторони алгоритмів за результатами експериментів}
\label{tab:strengths} \\
\hline
\textbf{Алгоритм} & \textbf{Сильні сторони} & \textbf{Слабкі сторони} \\
\hline
\endfirsthead
\hline
\textbf{Алгоритм} & \textbf{Сильні сторони} & \textbf{Слабкі сторони} \\
\hline
\endhead
Логістична регресія & найменше перенавчання (розрив ROC-AUC 0,010); відкалібровані ймовірності (Брієр 0,079); найшвидше навчання на підготовлених даних (1,4~с); найстійкіша до аугментації (втрата 0,009--0,011 PR-AUC); коефіцієнти можна інтерпретувати як внесок ознак у шанс повторної госпіталізації & найнижча якість ранжування (PR-AUC 0,172 проти 0,177--0,180): лінійна модель не бачить взаємодій ознак; найбільше залежить від обробки даних (без масштабування й кодування PR-AUC 0,166, навчання в 21 раз повільніше) \\
\hline
Випадковий ліс & найкраща якість без обробки даних (на тестовій вибірці PR-AUC 0,190 у сценарії А); не потребує масштабування; стійкий до вибору гіперпараметрів (за листка 3--30 PR-AUC змінюється не більш як на 0,002); дерева навчаються паралельно & запам'ятовує навчальні дані (ROC-AUC 0,985, а за листка 1 --- 1,000), тому потребує обмеження розміру листка; ймовірності зсунуті до середини, за порога 0,5 не позначає жодного пацієнта; чутливий до аугментації (втрата до 0,046 PR-AUC); найдовше навчання на підготовлених даних (2,5~с) і найбільший обсяг моделі в пам'яті \\
\hline
Градієнтний бустинг & найвища якість ранжування на крос-валідації (PR-AUC 0,179--0,180); швидке навчання завдяки групуванню значень в інтервали (1,7~с); відкалібровані ймовірності (Брієр 0,079); працює з ознаками різної природи без перетворень & найбільше страждає від аугментації: ROC-AUC з SMOTE 0,537, близько до випадкового вгадування, з великим розкидом між блоками; розрив ROC-AUC 0,215 без аугментації; найбільше гіперпараметрів (кількість дерев, крок навчання, розмір дерева, регуляризація) \\
\hline
\end{longtable}
}

Вибір алгоритму залежить від вимог до моделі. Якщо головне --- якість ранжування пацієнтів за ризиком, слід обрати бустинг або випадковий ліс: їхня перевага над логістичною регресією стабільна, а між собою вони рівні в межах розкиду. Якщо потрібні пояснення прогнозу чи робота на погано підготовлених даних із синтетичними записами, логістична регресія надійніша, ціною 0,005--0,012 PR-AUC.

\section{Перевірка на тестовій вибірці}
\label{sec:test}

Кожну пару «сценарій --- модель» навчено на всій навчальній вибірці та один раз застосовано до 13\,998 пацієнтів тестової вибірки. Поріг рішення взято з крос-валідації (табл.~\ref{tab:cv_tuned}). Результати наведено в табл.~\ref{tab:test} та на рис.~\ref{fig:test}.

\begin{table}[H]
\caption{Результати на тестовій вибірці за порога, обраного на крос-валідації}
\label{tab:test}
\footnotesize
\renewcommand{\arraystretch}{1.15}
\setlength{\tabcolsep}{4pt}
\begin{tabularx}{\linewidth}{|X|c|c|c|c|c|c|c|}
\hline
\textbf{Модель} & \textbf{ROC-AUC} & \textbf{PR-AUC} & \textbf{Поріг} & \textbf{Accuracy} & \textbf{Precision} & \textbf{Recall} & \textbf{F1} \\
\hline
\multicolumn{8}{|l|}{\textit{А: мінімальне очищення}} \\
\hline
логістична регресія & 0,649 & 0,168 & 0,121 & 0,783 & 0,172 & 0,369 & 0,234 \\
\hline
випадковий ліс & 0,663 & 0,190 & 0,120 & 0,797 & 0,187 & 0,374 & 0,249 \\
\hline
градієнтний бустинг & 0,652 & 0,186 & 0,105 & 0,724 & 0,152 & 0,455 & 0,228 \\
\hline
\multicolumn{8}{|l|}{\textit{Б: повне очищення}} \\
\hline
логістична регресія & 0,656 & 0,178 & 0,103 & 0,740 & 0,157 & 0,434 & 0,231 \\
\hline
випадковий ліс & 0,656 & 0,187 & 0,117 & 0,772 & 0,172 & 0,404 & 0,241 \\
\hline
градієнтний бустинг & 0,654 & 0,189 & 0,134 & 0,832 & 0,203 & 0,297 & 0,241 \\
\hline
\multicolumn{8}{|l|}{\textit{В: очищення + SMOTE}} \\
\hline
логістична регресія & 0,643 & 0,168 & 0,587 & 0,781 & 0,163 & 0,348 & 0,222 \\
\hline
випадковий ліс & 0,614 & 0,133 & 0,296 & 0,719 & 0,136 & 0,399 & 0,203 \\
\hline
градієнтний бустинг & 0,516 & 0,101 & 0,140 & 0,167 & 0,091 & 0,926 & 0,166 \\
\hline
\multicolumn{8}{|l|}{\textit{В: очищення + ADASYN}} \\
\hline
логістична регресія & 0,642 & 0,167 & 0,601 & 0,809 & 0,174 & 0,302 & 0,221 \\
\hline
випадковий ліс & 0,599 & 0,124 & 0,263 & 0,529 & 0,112 & 0,616 & 0,190 \\
\hline
градієнтний бустинг & 0,560 & 0,115 & 0,257 & 0,303 & 0,097 & 0,811 & 0,173 \\
\hline
\end{tabularx}
\end{table}

\begin{figure}[H]
  \centering
  \includegraphics[width=\linewidth]{lab2/test_curves}
  \caption{Криві точність-повнота та ROC-криві на тестовій вибірці}
  \label{fig:test}
\end{figure}

Тестова вибірка підтвердила висновки крос-валідації:

\begin{itemize}
\item повне очищення покращило логістичну регресію (PR-AUC 0,178 проти 0,168) і не вплинуло на деревоподібні моделі: випадковий ліс 0,190 і 0,187, бустинг 0,186 і 0,189. Ці різниці менші за розкид між блоками на крос-валідації;
\item аугментація погіршила всі моделі: логістичну регресію на 0,010--0,011 PR-AUC, випадковий ліс на 0,054--0,063, бустинг на 0,074--0,088. ROC-AUC бустингу з SMOTE становить 0,516, тобто на тестовій вибірці він ранжує пацієнтів майже випадково;
\item найкращі результати мають випадковий ліс у сценарії А (PR-AUC 0,190, ROC-AUC 0,663, F1 0,249) та бустинг у сценарії Б (PR-AUC 0,189). Базовий рівень PR-AUC випадкового класифікатора становить 0,090.
\end{itemize}

Пороги, обрані на крос-валідації, добре переносяться на нові дані, якщо модель навчалася без аугментації. Бустинг у сценарії Б з порогом 0,134 відносить до групи ризику 1\,834 пацієнтів (13,1~\% тестової вибірки). Серед них 373 справді повернулися до лікарні протягом 30 днів: precision 0,203, у 2,3 раза вище за базову частоту, recall 0,297. F1-міра на тесті (0,241) навіть вища, ніж на крос-валідації (0,233). Натомість бустинг з SMOTE з порогом 0,140 відносить до групи ризику 12\,731 пацієнта (91~\% вибірки), і його precision 0,091 не відрізняється від базової частоти. Імовірності моделі, навченої на синтетичних даних, нестабільні, і поріг, обраний на крос-валідації, на нових даних не працює.

За порога 0,5 accuracy моделей сценаріїв А і Б на тестовій вибірці становить 0,910--0,911, як і в тривіального правила «завжди клас 0». За обраних порогів accuracy нижча (0,72--0,83), але моделі знаходять 30--46~\% усіх повторних госпіталізацій. Оцінка Брієра моделей без аугментації (0,0785--0,0795) нижча за рівень сталого прогнозу 0,082, а з аугментацією становить 0,109--0,227.

\section{Висновки}

У лабораторній роботі порівняно три алгоритми класифікації --- логістичну регресію, випадковий ліс і гістограмний градієнтний бустинг --- на чотирьох варіантах підготовки даних Diabetes 130-US Hospitals для задачі передбачення повторної госпіталізації протягом 30 днів.

\begin{enumerate}
\item Accuracy непридатна для цієї задачі: усі моделі без аугментації мають 0,910 за порога 0,5, рівно стільки ж, скільки правило «завжди клас 0», і майже нікого не відносять до класу 1. Коректне порівняння дають метрики, що не залежать від порога (ROC-AUC, PR-AUC), та F1-міра за порога, обраного на навчальних даних.

\item Повне очищення даних важливе для логістичної регресії: PR-AUC зросла з 0,166 до 0,172 на крос-валідації та з 0,168 до 0,178 на тестовій вибірці, розрив між навчальною і валідаційною якістю зменшився вп'ятеро, а навчання прискорилося в 21 раз. Деревоподібні моделі до обробки байдужі: різниця між сценаріями А і Б для них менша за розкид між блоками.

\item Найвищу якість ранжування мають деревоподібні ансамблі: бустинг (PR-AUC 0,179--0,180 на крос-валідації) і випадковий ліс (0,177--0,178, а на тестовій вибірці 0,190). Логістична регресія відстає на 0,005--0,012 PR-AUC, але має найменше перенавчання і найпростіше пояснення прогнозу. Усі моделі без аугментації дають відкалібровані ймовірності (Брієр 0,079).

\item Аугментація SMOTE і ADASYN не покращила жодної моделі. Видиме зростання recall і F1-міри логістичної регресії за порога 0,5 є лише зсувом ймовірностей: за правильно обраного порога F1-міра з аугментацією нижча (0,218 проти 0,229), а оцінка Брієра зростає втричі. Деревоподібні моделі втратили 0,041--0,067 PR-AUC на крос-валідації; ROC-AUC бустингу з SMOTE на тестовій вибірці становить 0,516.

\item Аугментація викликала перенавчання особливого роду: деревоподібні моделі відносять до класу 1 99,9--100~\% синтетичних записів і лише 1,8--10,2~\% реальних пацієнтів із повторною госпіталізацією, на яких навчалися. Вони навчилися відрізняти синтетичні записи від реальних за дробовими значеннями One-Hot ознак і проміжними значеннями Target Encoding. Регуляризація моделі цього не усуває, а збільшення частки синтетичних записів посилює.

\item Аугментація призводить до перенавчання, коли модель гнучка, синтетичні записи мають неможливі для реальних даних значення (дискретні та закодовані ознаки), їх частка велика, а кодування навчальних і нових записів різне. Розрив між навчальною та валідаційною якістю цього не виявляє; виявляє порівняння частки розпізнаних синтетичних записів із recall на реальних даних.
\end{enumerate}

Головний висновок роботи: вибір алгоритму та вибір підготовки даних не можна робити окремо. Логістична регресія потребує ретельної обробки, але стійка до синтетичних даних; деревоподібні ансамблі дають кращу якість без обробки, але найбільше страждають від аугментації. Для незбалансованої задачі з дискретними ознаками найкращим рішенням виявилося не синтезувати дані, а навчати модель на реальних даних і обирати поріг рішення на крос-валідації.

\section*{Список використаних джерел}

\begin{enumerate}[label={[\arabic*]},ref={\arabic*},leftmargin=*,labelindent=0pt]
\item \label{ref:strack} Strack B., DeShazo J. P., Gennings C. et al. Impact of HbA1c Measurement on Hospital Readmission Rates: Analysis of 70,000 Clinical Database Patient Records // BioMed Research International. 2014. Vol. 2014. Article ID 781670.
\item \label{ref:smote} Chawla N. V., Bowyer K. W., Hall L. O., Kegelmeyer W. P. SMOTE: Synthetic Minority Over-sampling Technique // Journal of Artificial Intelligence Research. 2002. Vol. 16. P. 321--357.
\item \label{ref:adasyn} He H., Bai Y., Garcia E. A., Li S. ADASYN: Adaptive Synthetic Sampling Approach for Imbalanced Learning // Proc. IEEE International Joint Conference on Neural Networks. 2008. P. 1322--1328.
\item \label{ref:target} Micci-Barreca D. A Preprocessing Scheme for High-Cardinality Categorical Attributes in Classification and Prediction Problems // ACM SIGKDD Explorations Newsletter. 2001. Vol. 3, No. 1. P. 27--32.
\item \label{ref:breiman} Breiman L. Random Forests // Machine Learning. 2001. Vol. 45, No. 1. P. 5--32.
\item \label{ref:friedman} Friedman J. H. Greedy Function Approximation: A Gradient Boosting Machine // The Annals of Statistics. 2001. Vol. 29, No. 5. P. 1189--1232.
\item \label{ref:lightgbm} Ke G., Meng Q., Finley T. et al. LightGBM: A Highly Efficient Gradient Boosting Decision Tree // Advances in Neural Information Processing Systems 30. 2017. P. 3146--3154.
\item \label{ref:libsvm} Chang C.-C., Lin C.-J. LIBSVM: A Library for Support Vector Machines // ACM Transactions on Intelligent Systems and Technology. 2011. Vol. 2, No. 3. Article 27.
\item \label{ref:saito} Saito T., Rehmsmeier M. The Precision-Recall Plot Is More Informative than the ROC Plot When Evaluating Binary Classifiers on Imbalanced Datasets // PLoS ONE. 2015. Vol. 10, No. 3. e0118432.
\item \label{ref:niculescu} Niculescu-Mizil A., Caruana R. Predicting Good Probabilities with Supervised Learning // Proc. 22nd International Conference on Machine Learning. 2005. P. 625--632.
\item \label{ref:elor} Elor Y., Averbuch-Elor H. To SMOTE, or not to SMOTE? // arXiv preprint arXiv:2201.08528. 2022.
\item \label{ref:goorbergh} van den Goorbergh R., van Smeden M., Timmerman D., Van Calster B. The harm of class imbalance corrections for risk prediction models: illustration and simulation using logistic regression // Journal of the American Medical Informatics Association. 2022. Vol. 29, No. 9. P. 1525--1534.
\end{enumerate}

\end{document}
